% !TEX root = ../main.tex
\chapter{Tangent spaces}\label{ch:c04}
\begin{paracol}{2}

\begin{sources}
\begin{tabularx}{\linewidth}{@{}l l L@{}}
\toprule
\textbf{Video} & \textbf{Length} & \textbf{Content}\\
\midrule
\seg{4}{1}{Lecture 4: Tangent spaces (Tangentialräume)} & 1:44:23 & Extrinsic versus intrinsic velocity (the car and the trees);
smooth functions and curves; pointwise addition and scaling, the vector space $C^\infty(M)$; the tangent vector as a
directional derivative along a curve; the tangent space $T_pM$ and its vector space structure; components in a chart, the
induced basis $\partial/\partial x^j$; the chain-rule computation; change of chart: transformation of the basis and of the
components; the dual space, the cotangent space $T_p^*M$, the differential $d_pf$ and the coordinate differentials $d_px^i$.\\
\bottomrule
\end{tabularx}

\smallskip
\textbf{Transcript} \file{materials/transcripts/C04-V01.txt} (German YouTube auto-captions).
\end{sources}

\begin{reading}[Where this lecture sits in the other notes]
\emph{Geometric Anatomy}: lecture 10 (construction of the tangent space: smooth curves, directional derivatives, charts,
change of basis) and lecture 8 (dual vector space, covectors).
\emph{Gravity and Light} (TP-1 version): lecture 5 (tangent spaces).
\end{reading}

% ==================================================================
\section{Intuition: extrinsic and intrinsic velocity}\label{sec:ts-intuition}

Today: tangent vectors, ``one of the technical key concepts of mechanics'', built on what has been done so far
\ts{4}{1}{0:00}. The intuition is familiar: a curve on a differentiable manifold, and at each moment the
tangent vector to the curve, which in physics is the \emph{velocity vector} \ts{4}{1}{0:22}.
One could embed the (say two-dimensional) manifold in a higher-dimensional space (three-dimensional) and then understand
the vectors sticking out of the manifold \ts{4}{1}{1:07}. We do it differently, for a technical and a
philosophical reason: if the manifold is the whole universe or spacetime, it would be odd to have to embed it into something
larger to speak of velocities; that would lie outside the universe \ts{4}{1}{1:29}.

\begin{definition}[New concepts: extrinsic and intrinsic]\label{def:ts-intrinsic}
In differential geometry a definition that uses an embedding into a surrounding space is called \textbf{extrinsic}; one that
makes no reference to the surrounding space is \textbf{intrinsic}. We want an intrinsic definition of the velocity vector
\ts{4}{1}{2:10}.
\end{definition}
\begin{intuition}[The child in the car]
(The picture on the new problem sheet, top right, is headed ``vector counter'' [?], a joke on the driver.) Your earliest
memories of velocity are from the car. Velocity is never just \emph{there}: it is attached to a curve, here the track of the
car. As a child you looked out of the window: trees and posts at the side of the road flowing past; you could also look
at a mountain range getting higher and lower, or at the temperature changing (hot, cold) as you drive: the faster you drive,
the faster these \emph{reference functions} change, for given mountains, trees and temperature distribution.
So the velocity is encoded in how fast the values of functions in your surroundings change along your trajectory. That is the
intrinsic description: no surrounding space is needed, only the manifold you drive on and functions on it
\ts{4}{1}{2:31}.
\end{intuition}

% ==================================================================
\section{Smooth functions and smooth curves}\label{sec:ts-smooth}

Recall: a smooth manifold is a set $M$ with a topology and a $C^\infty$ atlas \ts{4}{1}{5:23}. A function
$f\colon M\to\R$ is \textbf{smooth} (we also say $C^\infty$) if for every chart $(U,x)$ of the atlas the chart
representation $f\circ x^{-1}\colon x(U)\subseteq\R^d\to\R$ is a smooth function (arbitrarily often continuously
differentiable), which one can check with elementary analysis \ts{4}{1}{6:04}.
(To remember the conditions Schuller does not memorise the formulas but redraws the picture: chart domain $U\subseteq M$,
chart map $x\colon U\to x(U)\subseteq\R^d$, and $f\circ x^{-1}$ as the representation of the function in the chart instead
of in the real world.)

The same for a curve $\gamma\colon I\subseteq\R\to M$: it is \textbf{smooth} if for every chart $(U,x)$ the chart
representation $x\circ\gamma\colon\gamma^{-1}(U)\subseteq\R\to x(U)\subseteq\R^d$ is smooth as a map between subsets of
$\R$ and $\R^d$ (to be exact, $\gamma$ restricted to the preimage of $U$ is a different map with a different domain, but
defined in the same way) \ts{4}{1}{7:47}.

\subsection{Pointwise operations}\label{ssec:ts-pointwise}

In school physics one adds curves, but that does not make sense: for curves $\gamma,\tilde\gamma$ one would add
pointwise, and in the target there is no addition: ``what is Paris plus Rome? it does not exist'' \ts{4}{1}{10:32}.
In a chart one could define the sum of the chart representations, but the result would not be chart independent: curves
cannot be added. \emph{Functions}, however, can \ts{4}{1}{10:53}.

\begin{definition}[New concepts: pointwise addition and scaling, pointwise inheritance]\label{def:ts-pointwise}
For functions $f,g\colon M\to\R$ and $\alpha\in\R$ define the new functions $f\oplus g$ and $\alpha\odot f$ by
\[
  (f\oplus g)(p):=f(p)+g(p),\qquad (\alpha\odot f)(p):=\alpha\cdot f(p)\qquad\text{for all }p\in M,
\]
where on the right-hand sides one has the known addition and multiplication on $\R$. This construction principle is called
\textbf{pointwise inheritance} (punktweise Vererbung): we know how to add in the target and do it point by point
\ts{4}{1}{11:39}.
\end{definition}
(Later the circles around $+$ and $\cdot$ are dropped; they are shown here once to make clear that this addition of
functions is not the addition of numbers.) \ts{4}{1}{13:03}
This is not unnecessary abstraction: the sum rule for derivatives at school refers to exactly this $+$; it differentiates
\emph{functions}, not function values \ts{4}{1}{15:33}.

\begin{definition}[New concepts: the set $C^\infty(M)$]\label{def:ts-cinfty}
$C^\infty(M)$ denotes the set of all smooth functions $M\to\R$. With pointwise addition and scaling it satisfies the axioms
of a \textbf{vector space} (commutativity, associativity, neutral element, inverse element, associativity of scaling, two
distributive laws, and the neutral scalar $1$) \ts{4}{1}{16:14}.
\end{definition}
Are these already the vectors we are looking for? No: they have nothing to do with the vectors we want to build, but we need
them as a tool. Do not let the word ``vector'' be taken too early \ts{4}{1}{17:59}.
\begin{exercise}[$C^\infty(M)$ is a vector space]\label{ex:ts-cinfty-vs}
Check that the sum and the scaling of smooth functions are smooth, and the eight vector space axioms
(it ``would have to be checked'' \ts{4}{1}{16:55}).
\end{exercise}

% ==================================================================
\section{Tangent vectors}\label{sec:ts-tangent}

Now the central point \ts{4}{1}{18:20}. Consider a smooth curve $\gamma\colon I\to M$ that passes
through a point $p\in M$ at parameter value $\lambda_0$: $\gamma(\lambda_0)=p$ \ts{4}{1}{19:02}.

\begin{definition}[New concepts: tangent vector, directional derivative]\label{def:ts-tangent}
The \textbf{tangent vector} to the curve $\gamma$ at the point $p=\gamma(\lambda_0)$ is the map
\[
  X_{\gamma,p}\colon C^\infty(M)\to\R,\qquad
  X_{\gamma,p}(f):=(f\circ\gamma)'(\lambda_0),
\]
that is: restrict $f$ to the curve ($f\circ\gamma\colon I\to\R$ is a function of the parameter alone), differentiate with
respect to the parameter, and evaluate at $\lambda_0$. It is a \emph{directional derivative}
(Richtungsableitung) of $f$ along $\gamma$ at $p$
\ts{4}{1}{19:47}.
\end{definition}
It is a linear map: $X_{\gamma,p}(f+g)=X_{\gamma,p}(f)+X_{\gamma,p}(g)$ and $X_{\gamma,p}(\alpha f)=\alpha X_{\gamma,p}(f)$, since
the sum and scaling of functions are pointwise and the derivative is linear; the domain and the target are both vector
spaces, so ``linear'' makes sense \ts{4}{1}{22:32}. (It also obeys a product rule.)
\ts{4}{1}{23:56}

\begin{intuition}[What you saw as a child]
This is what you perceived as a child as velocity; school deformed it into vector components. You sit in a car and drive
through a forest of trees of different heights (the function $f$): you do not look at how the trees change far away, only
\emph{where you are} (the composition $f\circ\gamma$): five metres high, two metres, two and a half, four; how does the
height of the trees I see change as I go? If you run the \emph{same path faster} (a different curve with the same trace),
the heights change faster. At the moment $\lambda_0$ the rate of change of the height of the trees where I am is the
directional derivative. From the same point in another direction it would be different, so this encodes the idea of the
direction \ts{4}{1}{24:17}.
\end{intuition}

\paragraph{Terminology and a warning.} One calls $X_{\gamma,p}$ the \textbf{tangent vector to the curve $\gamma$ at the point
$p$} (``in the point $p$ that the curve passes through''). One may think of an arrow, but from now on one must think of
a \emph{map}: a tangent vector is a directional derivative; ``in the differential geometry that we do, the concept of a
vector survives in the shape of a directional derivative''. One usually thinks the directional derivative is defined by a
direction; in reality the directional derivative is the elementary concept and the direction is the philosophy that goes
with it \ts{4}{1}{26:55}. Examples on
problem sheet~4, exercise~1 \ts{4}{1}{28:59}.

\begin{exercise}[Tangent vectors, problem sheet 4, exercise 1]\label{ex:ts-sheet4}
Examples of tangent vectors to curves on concrete manifolds, applied to concrete functions (new problem sheet, exercise~1).
\end{exercise}

Where do the velocity components of school come from? From the choice of a chart; they depend on the chart; this
comes below \ts{4}{1}{29:26}.

\subsection{The tangent space}\label{ssec:ts-space}

\begin{definition}[New concepts: tangent space]\label{def:ts-space}
The \textbf{tangent space} to $M$ at $p$ is the set
\[
  T_pM:=\{X_{\gamma,p}\mid \gamma\text{ a smooth curve through }p\},
\]
the collection of all tangent vectors at $p$ to all smooth curves through $p$ \ts{4}{1}{30:08}.
\end{definition}
Picture (on a patch of a sphere): through the point $p$ run many curves $\gamma$, also curves with the same trace but traversed
twice as fast, whose tangent vector is ``twice as long'': if one \emph{drew} the vectors as arrows (which one does not do,
since it is not an arrow but a map), the arrow of the faster curve would be twice as long; with all curves through $p$ and
all speeds the endpoints would fill the whole tangent plane \ts{4}{1}{30:50}.
We do not draw arrows because this is not the extrinsic picture but the intrinsic one: the price for not referring to the
surrounding space is that vectors are no longer arrows but directional derivatives, which act on
\emph{every} function (each ``temperature distribution'') defined on the manifold \ts{4}{1}{33:17}.

\paragraph{A student question.} (Garbled in the captions) A question on large velocities and on tanks and the sphere
\ts{4}{1}{34:38}. Schuller's answer: velocity is a property of a curve; ask ``what velocities
exist at this point in space, say my fingertip'': they are \emph{all} the velocities you can realise with curves through
the point (one flies through with a rocket one way, another another way); physically this totality is the tangent space,
technically it is the set of directional derivatives along these curves acting on arbitrary functions
\ts{4}{1}{35:38}. ``All my words are of course only pretty talk: the tangent vector is
this map; you will get used to it'' \ts{4}{1}{36:18}.

\subsection{Vector space structure}\label{ssec:ts-vs}

The tangent space is just a set so far; the sum and scaling are again pointwise. Vectors at \emph{different} points are in
different sets, with which we have nothing to do \ts{4}{1}{37:00}. Define on $T_pM$
\[
  (X\oplus Y)(f):=X(f)+Y(f),\qquad (\alpha\odot X)(f):=\alpha\cdot X(f)\qquad(X,Y\in T_pM,\ f\in C^\infty(M)),
\]
``punktweise'' in the sense that the map is added by adding the real numbers it produces on each argument
$f$ \ts{4}{1}{37:21}.
With these operations $T_pM$ is a \textbf{vector space} (the tangent vector space at $p$): the vector addition of
velocities at the same point \ts{4}{1}{40:32}.
(Showing that the sum is again of the form $X_{\gamma,p}$, i.e.\ that the set is closed under these operations, is
the content of the statement; it goes through a chart, see below.) The circles are dropped from now on
\ts{4}{1}{41:34}.
\begin{exercise}[$T_pM$ is a vector space]\label{ex:ts-vs}
Show that $T_pM$ is closed under $\oplus$ and $\odot$ (find a curve whose tangent vector is $X_{\gamma,p}\oplus X_{\tilde\gamma,p}$,
e.g.\ using a chart) and satisfies the vector space axioms. (Said in class: ``one can show''.)
\end{exercise}

% ==================================================================
\section{Components in a chart}\label{sec:ts-components}

So far everything is abstract, ``and probably very different from what you thought about velocity, as a child''. Now the
contact with what you know: that in a chart velocity is a collection of components \ts{4}{1}{42:15}.
We introduce a chart \ts{4}{1}{43:18}.

Ingredients: a smooth curve $\gamma$ on $M$, an arbitrary smooth function $f$, and a chart $(U,x)$ of the smooth atlas
with $p\in U$ \ts{4}{1}{44:04}. The chart gives the chart
representation of the curve, $x\circ\gamma$ (the course the captain drew on the sea chart) and of the function,
$f\circ x^{-1}$ (the temperature distribution drawn on the sea chart) \ts{4}{1}{45:46}.

Compute (``most calculations are done by writing a chain of equalities starting from the left'') \ts{4}{1}{46:57}.
Using the definition, and inserting an identity $x^{-1}\circ x$ (it changes nothing) and the associativity of composition:
\begin{align*}
  X_{\gamma,p}(f)&=(f\circ\gamma)'(\lambda_0)=\bigl((f\circ x^{-1})\circ(x\circ\gamma)\bigr)'(\lambda_0)\\
  &=\sum_{j=1}^d\,\partial_j\bigl(f\circ x^{-1}\bigr)\bigl(x(p)\bigr)\cdot\bigl(x^j\circ\gamma\bigr)'(\lambda_0)
\end{align*}
\ts{4}{1}{47:38}.
The last step is the multidimensional chain rule for the composition $\R\xrightarrow{x\circ\gamma}\R^d\xrightarrow{f\circ x^{-1}}\R$: the outer derivative is the
derivative in the $j$-th argument (written $\partial_j$, the derivative with respect to the $j$-th slot, evaluated at
the point $x(\gamma(\lambda_0))=x(p)$); the inner derivative is the derivative of the $j$-th component of the curve
representation; one sums over $j$ (Einstein summation convention of problem sheet~1)
\ts{4}{1}{50:23}.

\begin{definition}[New concepts: the basis $\partial/\partial x^j$, components of a tangent vector]\label{def:ts-basis}
Write $\gamma^j_{(x)}:=x^j\circ\gamma$ for the $j$-th component of the chart representation of $\gamma$ (in class $q^j$),
and define the maps
\[
  \left(\frac{\partial}{\partial x^j}\right)_{\!p}\colon C^\infty(M)\to\R,\qquad
  \left(\frac{\partial}{\partial x^j}\right)_{\!p}f:=\partial_j\bigl(f\circ x^{-1}\bigr)\bigl(x(p)\bigr).
\]
Then
\[
  X_{\gamma,p}\;=\;\dot\gamma^{\,j}_{(x)}(\lambda_0)\left(\frac{\partial}{\partial x^j}\right)_{\!p},
  \qquad \dot\gamma^{\,j}_{(x)}(\lambda_0):=\bigl(x^j\circ\gamma\bigr)'(\lambda_0).
\]
The numbers $\dot\gamma^j_{(x)}(\lambda_0)$ are the \textbf{components of the tangent vector $X_{\gamma,p}$ with respect to the chart
$(U,x)$}, and $\{(\partial/\partial x^j)_p\}_{j=1,\dots,d}$ is the \textbf{chart-induced basis} of $T_pM$
\ts{4}{1}{54:59}.
\end{definition}
Remarks:
\begin{itemize}
  \item The symbol $\partial/\partial x^j|_p$ is a \emph{new symbol}, not simply ``$\partial f/\partial x^j$''. Schuller had asked
        on the first problem sheet never to write $\partial f/\partial x^j$ for the partial derivative; the partial derivative
        is the derivative with respect to the $j$-th slot of the function on $\R^d$. When you see the symbol, expand
        it as above (it contains the function, the chart map, the direction, and the point), then you may compute with the
        rules of differentiation on $\R^d$ \ts{4}{1}{55:40}.
  \item Later $f$ is dropped from the notation and the symbol is regarded as an operator \ts{4}{1}{57:33}.
  \item The components are \emph{numbers} (``$2,3,7$''): without a basis they mean nothing. A basis (in the sense of linear
        algebra: linearly independent generating system) is there once a chart is chosen; there are no unit vectors
        without a metric \ts{4}{1}{1:00:40}.
  \item The basis is ``distinguished'' only by the chart; with another chart one gets another distinguished basis, and the
        abstract object ``velocity vector in the real world'' cannot depend on the chart: so the \emph{components must change}
        when the basis changes, as in linear algebra \ts{4}{1}{1:02:48}.
\end{itemize}
In physics, with a chart chosen, the velocity has concrete components, obtained by differentiating the components of
the curve in the chart with respect to the parameter: that is exactly what school does \ts{4}{1}{1:03:30}.
But the lesson is that these numbers are not yet the velocity vector: they are only its components, meaningless unless one knows
with respect to which basis. By one small harmless calculation we identified at once what the abstract vector in the real
world is, what its components are, and with respect to which basis \ts{4}{1}{1:04:11}.
The importance of $T_pM\ni X\mapsto$ components is:
\begin{aside}[Generalised coordinates]
In standard mechanics books ``generalised coordinates and generalised velocities'' appear around the tenth lecture; here we
are in the fourth, because the conceptual material is brought forward to know what we are talking about
\ts{4}{1}{1:32:35}.
\end{aside}

% ==================================================================
\section{Change of chart}\label{sec:ts-chartchange}

After a break of ten minutes: what happens to the basis and the components if the chart is changed
\ts{4}{1}{1:04:53}? Situation: $X_{\gamma,p}$ is the abstract object (a directional
derivative); with respect to the chart $x$ we derived
\[
  X_{\gamma,p}=\dot\gamma^{\,j}_{(x)}(\lambda_0)\left(\frac{\partial}{\partial x^j}\right)_{\!p},
\]
and with respect to another chart $(V,y)$ the same abstract object reads
$X_{\gamma,p}=\dot\gamma^{\,j}_{(y)}(\lambda_0)(\partial/\partial y^j)_p$ (the components of the curve in the $y$-coordinates,
differentiated at the right parameter); ``this is two times the same abstract object, so the left and right sides must be the
same'' \ts{4}{1}{1:05:40}.
The practical question: given the components in one chart and the chart transition map, find the components in the other
chart \ts{4}{1}{1:09:29} (e.g.\ vectors given in Cartesian coordinates when the problem
is better solved in polar coordinates) \ts{4}{1}{1:10:10}.

Strategy: ``components times basis equals components times basis''; if both sides had the \emph{same basis} we could compare
components: two vectors with the same components with respect to the same basis are equal, because a basis is a linearly
independent generating system, and linear independence is exactly what allows one to conclude that the coefficients agree
\ts{4}{1}{1:10:30}.
So first express the basis $(\partial/\partial x^j)_p$ through $(\partial/\partial y^m)_p$.

\begin{theorem}[Transformation of the induced basis]\label{thm:ts-basis-trafo}
Let $(U,x)$ and $(V,y)$ be charts of the smooth atlas with $p\in U\cap V$. Then
\[
  \left(\frac{\partial}{\partial x^j}\right)_{\!p}
  \;=\;\partial_j\bigl(y^m\circ x^{-1}\bigr)\bigl(x(p)\bigr)\left(\frac{\partial}{\partial y^m}\right)_{\!p}.
\]
\end{theorem}
\begin{proof}
Apply the left side to an arbitrary $f\in C^\infty(M)$; by definition
\[
  \left(\frac{\partial}{\partial x^j}\right)_{\!p}f=\partial_j\bigl(f\circ x^{-1}\bigr)\bigl(x(p)\bigr).
\]
Insert $y^{-1}\circ y$ (a ``conjurer's trick'' that changes nothing): $f\circ x^{-1}=(f\circ y^{-1})\circ(y\circ x^{-1})$.
The two factors are: $f\circ y^{-1}\colon y(V)\subseteq\R^d\to\R$ (the chart representation of $f$ in the chart $y$) and
$y\circ x^{-1}\colon x(U\cap V)\to y(U\cap V)$ (the transition map, which is also needed in the picture)
\ts{4}{1}{1:11:55}.
By the multidimensional chain rule (outer derivative first, in the $m$-th slot of $f\circ y^{-1}$, evaluated at
$y(x^{-1}(x(p)))=y(p)$; the inner derivative is the derivative of the $m$-th component of the transition map in the $j$-th
slot):
\[
  \partial_j\bigl(f\circ x^{-1}\bigr)\bigl(x(p)\bigr)
  =\partial_m\bigl(f\circ y^{-1}\bigr)\bigl(y(p)\bigr)\cdot\partial_j\bigl(y^m\circ x^{-1}\bigr)\bigl(x(p)\bigr)
\]
(with a sum over $m$: an index appearing once up and once down is summed over; Einstein convention)
\ts{4}{1}{1:15:22}.
The first factor is $(\partial/\partial y^m)_pf$, so the claim follows for all $f$ \ts{4}{1}{1:20:18}.
\end{proof}
\begin{aside}[On the chain rule]
The multidimensional chain rule: ``what acts on everything stands at the inner derivative, and the rest is summed over the
repeated index''. If it is not familiar, look in a first-year engineering mathematics book (Schuller recommends
Papula, \emph{Mathematik für Ingenieure und Naturwissenschaftler}, with ``1000 examples''; probably vol.~2). This is ``first
semester material for physicists'' and the whole proof is only an exercise in the chain rule plus the definition of the
symbol \ts{4}{1}{1:17:47}.
\end{aside}

Now compare the two expressions of the \emph{same} vector (``we start at the right side and end on the left'')
\ts{4}{1}{1:23:40}: insert the theorem into
$\dot\gamma^j_{(x)}(\partial/\partial x^j)_p$:
\[
  X_{\gamma,p}=\dot\gamma^{\,j}_{(x)}(\lambda_0)\,\partial_j\bigl(y^m\circ x^{-1}\bigr)\bigl(x(p)\bigr)\left(\frac{\partial}{\partial y^m}\right)_{\!p}
  \;\overset{!}{=}\;\dot\gamma^{\,m}_{(y)}(\lambda_0)\left(\frac{\partial}{\partial y^m}\right)_{\!p}.
\]
Subtract the right from the left side (renaming the dummy index) \ts{4}{1}{1:24:22}:
\[
  0=\Bigl(\dot\gamma^{\,j}_{(x)}(\lambda_0)\,\partial_j\bigl(y^m\circ x^{-1}\bigr)\bigl(x(p)\bigr)-\dot\gamma^{\,m}_{(y)}(\lambda_0)\Bigr)\left(\frac{\partial}{\partial y^m}\right)_{\!p}.
\]
Because $\{(\partial/\partial y^m)_p\}$ is a basis, in particular linearly independent, every bracket must vanish
\ts{4}{1}{1:27:33}:

\begin{theorem}[Transformation of the components of a tangent vector]\label{thm:ts-comp-trafo}
\[
  \dot\gamma^{\,m}_{(y)}(\lambda_0)\;=\;\partial_j\bigl(y^m\circ x^{-1}\bigr)\bigl(x(p)\bigr)\;\dot\gamma^{\,j}_{(x)}(\lambda_0)
  \qquad(m=1,\dots,d),
\]
i.e.\ the components in the new chart are obtained from the components in the old chart by the Jacobian matrix of the chart
transition map $y\circ x^{-1}$ at the point $x(p)$ \ts{4}{1}{1:28:36}.
\end{theorem}
Interpretation: if someone gives you the velocity ``2, 3, 7'' with respect to some chart, and you change to another chart
via a transition map, then the $m$-th new component depends on \emph{all} old components through a linear map at this point,
a matrix (``a matrix is a linear map, of course'' with two indices) \ts{4}{1}{1:30:28}.

\begin{coursenote}[Index conventions]
Every index appearing twice (once up, once down) is summed over; the name of a summed index does not matter, but one must
avoid conflicts: no more than two indices of the same name in a product; when combining pieces one renames summation indices to
unused letters \ts{4}{1}{1:31:53}.
\end{coursenote}
This is ``what standard mechanics books call the change of generalised velocities under a change of generalised
coordinates''. It is essential and must be understood: ``if you understand this, nobody can fool you any more''
\ts{4}{1}{1:32:35}. After a short announcement (students are encouraged to come to the
office after the lecture or on Tuesday with questions about physics or mathematics) \ts{4}{1}{1:33:38}.

% ==================================================================
\section{The cotangent space}\label{sec:ts-cotangent}

One last concept of utmost importance, with exercises to follow \ts{4}{1}{1:34:31}.

\begin{definition}[New concepts: dual vector space, cotangent space]\label{def:ts-dual}
For a vector space $V$ over $\R$, the \textbf{dual space} (Dualraum) $V^*$ is the set of all linear maps $V\to\R$
(linear: $\varphi(v+w)=\varphi(v)+\varphi(w)$ and $\varphi(\alpha v)=\alpha\varphi(v)$; one first adds then maps, or maps then adds);
with the pointwise operations inherited from $\R$ it is again a vector space, of the same dimension if $V$ is
finite-dimensional \ts{4}{1}{1:34:55}.
The \textbf{cotangent space} (Kotangentialraum) at $p$ is the dual of the tangent space,
\[
  T_p^*M:=(T_pM)^*,
\]
the set of linear maps $T_pM\to\R$; the star is written in place of the brackets, ``because we are lazy''
\ts{4}{1}{1:37:20}.
\end{definition}
\begin{coursenote}[Dual space as ``companion'']
In German the dual space is also called the ``Dualraum'' or ``Zuweg-torraum'' [?: the captions garble it]; the dual vector
space is the companion of $V$, which is why one speaks of ``dual'' \ts{4}{1}{1:36:59}.
\end{coursenote}
A concrete example, known from the upper school:

\begin{definition}[New concepts: differential of a function at a point]\label{def:ts-differential}
For $p\in M$ and $f\in C^\infty(M)$ define $d_pf\colon T_pM\to\R$ by
\[
  (d_pf)(X):=X(f)\qquad(X\in T_pM)
\]
\ts{4}{1}{1:38:23}.
\end{definition}
(Schuller first writes $d_pf$ on the board without the subscript, then notes it is defined only at the point $p$; he notices
that ``$d$ eats functions and gives elements of the cotangent space, not a real number''; the real number appears only
after feeding a tangent vector.) This is linear in $X$ because $(X+Y)(f)=X(f)+Y(f)$ and $(\alpha X)(f)=\alpha X(f)$ by the
pointwise definition, hence $d_pf\in T_p^*M$ \ts{4}{1}{1:40:13}.

In particular, for a chart $(U,x)$ the coordinate functions $x^1,\dots,x^d\colon U\to\R$ are smooth functions, so one may form
$d_px^1,\dots,d_px^d$. (A chart map is a priori only continuous; but by the definition of smoothness the $x^i$ are smooth
functions on $U$, ``try to show this as a homework'' \ts{4}{1}{1:41:45}.) These are the ``$dx$'' of
school: in dimension one the $dx$ in the integral, here the formal definition; in higher dimension they are also the $dx^i$
that occur in integrals \ts{4}{1}{1:42:06}.
\begin{exercise}[Chart coordinate functions are smooth; the gradient is a covector]\label{ex:ts-grad}
Show that the coordinate functions $x^i$ of a chart in a smooth atlas are smooth functions on the chart domain. On the
next problem sheet: the \emph{gradient of a function is not a vector but a covector} (``the gradient is a covector and not
a vector, one has the habit to confuse them''), to be practised there \ts{4}{1}{1:43:10}.
\end{exercise}

\paragraph{A student remark: $C^1$ suffices.} A student notes that $C^1$ would suffice. Schuller: right, I should have said
so; but there are theorems that guarantee that from any $C^1$ atlas one can make a $C^\infty$ atlas (roughly speaking), so one
simplifies the assumption and does not care how often one can differentiate; ``one could express it more sharply, you are
completely right'' \ts{4}{1}{1:43:30}.

% ==================================================================
\section*{Key concepts}
\begin{keyconcepts}
\textbf{Retrieval questions.}
\begin{enumerate}
  \item What is the difference between an extrinsic and an intrinsic definition of a tangent vector, and why do we want the second?
  \item Why can functions be added pointwise but curves not? What does ``pointwise inheritance'' mean?
  \item Define $X_{\gamma,p}$ and show it is linear. What kind of object is a tangent vector?
  \item Define $T_pM$, its vector space operations; why do vectors at different points not add?
  \item Derive the components $\dot\gamma^j$ and the basis $\partial/\partial x^j|_p$ from the chain rule. Why does a basis
        need a chart?
  \item Derive the transformation of the basis and of the components under a change of chart.
  \item Define $V^*$, $T_p^*M$, $d_pf$ and $d_px^i$. Why is the gradient a covector?
\end{enumerate}
\textbf{Concepts.}
\begin{itemize}
  \item Extrinsic vs.\ intrinsic (\cref{def:ts-intrinsic}); $C^\infty(M)$ as a vector space (\cref{def:ts-cinfty}).
  \item Tangent vector as directional derivative along a curve (\cref{def:ts-tangent}); tangent space (\cref{def:ts-space}).
  \item Components and induced basis (\cref{def:ts-basis}); transformation rules
        (\cref{thm:ts-basis-trafo,thm:ts-comp-trafo}).
  \item Dual space, cotangent space, differential $d_pf$ (\cref{def:ts-dual,def:ts-differential}).
\end{itemize}
\end{keyconcepts}

\flushside
\end{paracol}
